\section{Structure of one-loop corrections}

\subsection{Confinement and infrared cut-offs}

 There are several standard techniques to calculate gluon radiative corrections in QCD.
 Most of them  are oriented to work in the region of
 absolute perturbative QCD (pQCD) applicability.
A
  straightforward use of such methods, however,
 may need some  care  in applications involving energy scales
 that are not very large.
   For this reason,  let us discuss
   some features of calculations
   on the border of applicability of perturbative methods.


To begin with, one should  remember that quarks and gluons are confined, i.e.
 the  propagators of all   diagrams  (even in a continuum case)
 are embedded in a finite volume whose size is determined by the hadron's radius.
 The confinement  effects lead, in particular,
  to a rapid   decrease of  correlators like ITDs or pseudo-PDFs  at  distances $z_3$
larger  than the  hadronic radius  $R$.
Still,  at short distances  one can  use  asymptotic freedom
and obtain, in particular, the $\ln z_3^2$ singularities.

Thus, it makes sense to treat  pseudo-ITDs and  pseudo-PDFs  as  sums  of the soft and hard parts.
The soft  part  basically
reflects  the  size of the system  and is assumed to be  finite for $z_3=0$.   The hard part
is   singular for \mbox{$z_3 \to 0$},   and  is   produced by perturbative interactions.
The hard part may be visualized  then as
generated from the soft part  through a hard exchange kernel $ H (0,z; z_1, z_2)$,
\begin{align}
    & {\cal M}^{\rm   hard} (\nu, -z^2=z_3^2)
    \nn & =   \int d^4z_1 \, d^4 z_2 \, H (0,z; z_1, z_2) \,  {\cal M}^{\rm  soft} (z_1,z_2)\  .
\label{Mhar}
\end{align}

In the standard  pQCD factorization approaches,
the soft part is mimicked  by on-shell parton states, and the
 \mbox{$\ln z_3^2$-singularities}  appear either as $\ln (z_3^2 m^2)$,
where $m$ is the parton mass or  $\ln (z_3^2 \mu_{\rm IR}^2)$,
where $\mu_{\rm IR}$ is the scale used in dimensional regularization
of infrared singularities in  the  case of massless partons.



Since   ${\cal  M}^{\rm  soft} (z_1,z_2)$ in Eq. (\ref{Mhar})
rapidly decreases for large separations
$|z_1-z_2|$,  the hadronic size $R$  provides an infrared cut-off for the integral,
even when the quarks are massless.  While at   short distances   one   gets  the
 $\ln (z_3^2 / R^2)$ behavior,   the logarithmic form  is just an approximation
valid for $z_3 \ll R$. Such
a restriction may be hard to implement on the lattice.


 Of course, the exact form of the IR regularization  imposed
by confinement is not known. To get  a feeling,
 let us take an infrared regularization by a  mass term.
 A typical integral producing the  $\ln z_3^2$  singularity then  has the form
   \begin{align}
  I_K (z_3^2) = \int_0^\infty \frac{d\alpha}{\alpha} e^{-z_3^2/4\alpha -   \alpha m^2} \ ,
  \label{IK}
     \end{align}
where $\alpha$ is the Schwinger's $\alpha$-parameter and $m$ is the infrared  regulator.
 One can see that
   \begin{align}
  I_K (z_3^2) = &2 K_0 (mz_3) \nn & = - \ln (m^2 z_3^2) + 2 \ln (2e^{-  \gamma_E })+ {\cal O} (z_3^2)  \  ,
  \label{Kexp}
     \end{align}
where $K_0(mz_3)$ is the modified Bessel function.  Its expansion for small $z_3$
explicitly  shows the expected $\ln (z_3^2 m^2)$ singularity.


The usual pQCD factorization   procedure   is to split
$\ln (z_3/ R)$ into the short-distance part $\ln (z_3 \mu)$   that is attributed  to
 the coefficient function  and the long-distance part $\ln ( 1/ \mu R )$
that is   absorbed into the ``renormalized'' PDF $f(x,\mu^2)$.
Given the commonly used   lattice   spacing $a\sim 0.1$ fm  and the hadron size
$R \lesssim 1$ fm, the question is whether  there is enough interval
for the logarithmic part
of the \mbox{$z_3$-dependence} to   be visible  in the  data  at all.

An important feature of the Bessel  function $K_0 (mz_3)$ is that it
  exponentially decreases when
$z_3$  exceeds the infrared cut-off $1/m$.
Thus, if   instead of the  short-distance approximation of
$I_K(z_3^2)$ by $\ln (1/z_3^2)$,
one would use the ``exact'' $I_K(z_3^2)$ function
for the evolution term,
 there will be no evolution corrections for large $z_3$.
In other words, the logarithmic
evolution disappears at large distances.



\subsection{Rescaling relation}

To fix a  relation between the pseudo-PDF scale $z_3$  and the $\overline{\rm MS}$ scale $\mu$,
one should take into account constant terms, like $2 \ln (2e^{-  \gamma_E })$  in Eq. (\ref{Kexp}).
In the $\overline{\rm MS}$-OPE approach,
  one takes $z^2=0$
and  then  applies
 the dimensional regularization which   adds  the  $\alpha^{\epsilon}$ factor into the integral
(\ref{IK}) making it convergent. After that, one    uses the $\overline{\rm MS}$-prescription,
which is   arranged  to produce  exactly
 $\ln (\mu^2/m^2)$ as the result in this case,
   \begin{align}
  I_{Dm} (\mu^2) =& \int_0^\infty \frac{d\alpha}{\alpha}  (\alpha \mu^2 e^{\gamma_E})^{\epsilon}  e^{-   \alpha m^2} \,
  =\Gamma (\epsilon) \left ( \frac{\mu^2 e^{\gamma_E}}{m^2} \right )^\epsilon \nn & \to
  \frac1{\epsilon}   + \ln (\mu^2/m^2)  \  .
  \label{ImM}
     \end{align}
 Thus,  the constant  term in Eq. (\ref{Kexp})
provides the leading-logarithm rescaling coefficient
$2e^{-  \gamma_E}$  between the pseudo-PDFs and $\overline{\rm MS}$ parton distributions
expressed by Eq. (\ref{Ptof}).





One may ask what happens  if one uses another type of    the IR regularization.
In particular,  the  Gaussian models for TMDs suggest that the
decrease for large  $z_3$ is also Gaussian.  One may expect
 that    the hard correction
should resemble,  for  large $z_3$,  the behavior of the soft part.
  Thus,   the exponential $e^{-m|z_3|}$ fall-off of the modified
Bessel function may  look   too slow.
A  Gaussian decrease can  be easily  provided by a sharp IR cut-off
    \begin{align}
  I_G (z_3^2) = \int_0^{z_0^2/4}  \frac{d\alpha}{\alpha}    e^{-z_3^2/4\alpha } = \Gamma [0,z_3^2/z_0^2]
  \label{IDexp}
     \end{align}
applied to  Eq. (\ref{IK}).  For small $z_3^2$,  the incomplete
gamma-function  $ \Gamma (0,z_3^2/z_3^2)$   has a logarithmic singularity
   \begin{align}
  \Gamma (0,z_3^2/z_0^2) =  \ln (z_0^2/z_3^2) -  \gamma_E + {\cal O} (z_3^2)  \  ,
  \label{Dexp0}
     \end{align}
     while for large $z_3^2$, the function  $I_G (z_3^2)$  has  a Gaussian  $e^{-z_3^2/z_0^2}$  fall-off.
Again, we  can  calculate the $z_3=0$ version of
Eq. (\ref{IDexp})  using the $\overline{\rm MS}$-scheme to obtain
  \begin{align}
  I_{DG} (\mu^2) = & \int_0^{z_0^2/4}  \frac{d\alpha}{\alpha}  (\alpha \mu^2  e^{\gamma_E})^{\epsilon}    =
  \frac1{\epsilon} \,  \left (\frac{ z_0^2  \mu^2  e^{\gamma_E}}{4}  \right )^{\epsilon} \nn &
  \to  \frac1{\epsilon} \, + \ln (z_0^2 \mu^2    ) -2  \ln ( 2 e^{-\gamma_E}) -  \gamma_E \ .
  \label{IDexp2}
     \end{align}
One can see that the pseudo-PDF/PDF rescaling (\ref{Ptof}) remains intact.
This   is a natural result, because the   relation between the finite-$z_3$
and $\overline{\rm MS}$  cut-offs  concerns only the short-distance  properties
of  the bilocal operator.


\subsection{One-loop correction}

The  discussion given in the  previous section   addresses  only the
overall  rescaling between  two regularization schemes
(just like    the  relation between the values of the QCD scale $\Lambda$ in, say, MOM and
$\overline{\rm MS}$ schemes).
To  establish a connection  between the pseudo-PDFs and  the
$\overline{\rm MS}$-PDFs, we  need, in addition, the constant part of the one-loop
coefficient function in the nonlocal OPE of  \mbox{Eq. (\ref{OPE}).  }
It   was given  in Refs. \cite{Ji:2017rah} and \cite{Radyushkin:2017lvu},
with some differences between them.
After rechecking our calculation and fixing  typos,
we present our result  in the form
  \begin{align}
  &
{\mathfrak M}  (\nu, z_3^2)    = {\mathfrak  M}^{\rm soft}  ( \nu, 0)  -  \frac{\alpha_s}{2\pi} \, C_F
\int_0^1  dw \,  \left \{   \frac{1+w^2} {1-w}    \right.  \nn & \left.  \hspace{2cm} \times
  \left [  \ln \left (z_3^2m^2\frac{  e^{2\gamma_E}}{4} \right )   +1 \right ]
\right. \nn    & \left.
+ 4  \,  \frac{\ln (1-w)}{1-w}
\right \}  \left [
{\mathfrak  M}^{\rm soft}  (w \nu,0) -  {\mathfrak  M}^{\rm soft}  ( \nu,0) \right ]\  .
 \label{Mhard}
 \end{align}



Turning to the PDF counterpart, we take $z^2=0$ and using the  $\overline{\rm MS}$ scheme for the UV divergence, obtain
   \begin{align}
{\cal I}  (\nu, \mu^2)
  =&  {\mathfrak  M}^{\rm soft}  ( \nu, 0) \nn & -  \frac{\alpha_s}{2\pi} \, C_F\
    \int_0^1  dw \
  \left [
{\mathfrak  M}^{\rm soft}  (w \nu, 0) -  {\mathfrak  M}^{\rm soft}  ( \nu, 0) \right ]\nn &
  \nn & \times \left \{    \frac{1+w^2} {1-w} \,  \ln (m^2/ \mu^2)   + 2 (1-w) \right \}
.
\label{Msbar}
 \end{align}

The logarithmic part here involves a convolution that  may be symbolically written as
   $
B \otimes {\mathfrak M}\,  (\nu) $ where
     \begin{align}
B(w) =
\left [ \frac{1+w^2}{1- w} \, \right ]_+
\label{AP}
 \end{align}
 is  the Altarelli-Parisi kernel  \cite{Altarelli:1977zs}.

Combining Eqs. (\ref {Mhard}) and (\ref{Msbar}), we obtain
 the relation
  \begin{align}
{\cal I}  (\nu, \mu^2)    = & {\mathfrak  M} ( \nu, z_3^2) +  \frac{\alpha_s}{2\pi} \, C_F\
 \int_0^1  dw \,   {\mathfrak  M}  (w \nu,z_3^2)   \nn &
 \times  \left  \{   B(w) \, \left [ \ln \left (z_3^2\mu^2 \frac{  e^{2\gamma_E}}{4} \right )  +1
 \right ]   \right. \nn & \left.
+   \left [ 4  \frac{\ln (1-w)}{1-w} - 2 (1-w)  \right ]_+ \right  \}
\
\label{MNL}
 \end{align}
 which is in agreement with a recent result of Ref. \cite{Izubuchi:2018srq}
 (see  also  Ref. \cite{Zhang:2018ggy}).
Eq. (\ref{MNL})   allows  one to convert the  data points
for ${\cal  M} ( \nu, z_3^2) $  into the ``data'' for ${\cal  I} (\nu, \mu^2)$.



The first contribution in   the   second line is  an obvious  term reflecting
 the  general multiplicative scale  difference between  the $z^2$ and $\overline{\rm MS}$  cut-offs.
If  all  the further  terms  are  neglected, then the only difference between
 $ {\mathfrak  M} ( \nu, z_3^2) $  and ${\cal  I} (\nu, \mu^2)  $
 is  just the rescaling \mbox{$\mu^2  =4e^{-2\gamma_E}/ z_3^2$.}
 In that case,   one can evolve the $ {\mathfrak  M} ( \nu, z_3^2) $  data to a particular
 $z_3$   value  $z_0$, and treat (in this approximation) the resulting function
 $ {\mathfrak  M} ( \nu, z_0^2) $ as the  $\overline{\rm MS}$  ITD
 corresponding to the scale \mbox{$\mu = 2 e^{-\gamma_E}/ z_0$,}
 which is numerically close to $1/z_0$.

 This   simple rescaling relation (used in Ref. \cite{Orginos:2017kos})    is modified
 when  the   further  terms  of Eq. (\ref{MNL})  are included.
 In particular,
 the  term  proportional
to the Altarelli-Parisi kernel $B(w)$ may be absorbed into the $\ln z_3^2$ term,
which would just  change the  rescaling relation
 into   $\mu = 2 e^{-1/2-\gamma_E}/ z_0$.




    \begin{figure}[tbp]
    \centerline{\includegraphics[width=3in]{images/linkboth}}
    \caption{Coordinate representation for  diagrams  producing a  large one-loop correction.
        \label{link}}
    \end{figure}


The term with $[\ln (1-w)]/(1-w)$  produces  a large negative contribution.
In Feynman gauge, according to Ref. \cite{Radyushkin:2017lvu},  it comes from the evolution part of the  vertex diagrams
  involving
    the gauge link      (see Fig.\ref{link}).  The key point is   that the gluon is attached there  to
    a running $tz_3$  position on the link.  After integration
    over $t$,  etc., the net outcome is  that the $z_3$-dependence
    of these diagrams  is generated by an effective  scale  smaller than $z_3$.
    Indeed, let  us   combine the  $[\ln (1-w)]/(1-w)$   term with the $\ln z_3^2$ logarithm
    by rewriting Eq. (\ref{MNL}) as
    \begin{align}
{\cal I}  (\nu, \mu^2)    = & {\mathfrak  M} ( \nu, z_3^2) +  \frac{\alpha_s}{\pi} \, C_F\
 \int_0^1  dw \,   {\mathfrak  M}  (w \nu,z_3^2)   \nn &
 \times  \Biggl   \{   \frac{1+w^2}{1-w}  \,  \ln \left [(1-w)z_3  \mu \frac{  e^{\gamma_E+1/2}}{2} \right ]
\nn &
+   \left [   (w+1)  \ln (1-w)  -(1-w) \right ] \Biggr  \} _+
\   .
\label{MNL2}
 \end{align}
 We  see that $z_3$  enters now through  a running $(1-w)z_3$  location.
 The remaining  $(w+1) \ln (1-w)$  term is much less  singular
 than $B(w)$ for $w=1$, and does not produce large contributions.

 Thus, the magnitude
 of the one-loop correction is governed by the combined evolution logarithm.
 It cannot be made zero by a particular choice of $\mu$  because
 it depends on the integration variable $w$.
 Still,   the $w$-integrated contribution will vanish for some $\mu$   that
 we may write as
     \begin{align}
\mu =  \frac{2 e^{-1/2-\gamma_E}}{\langle 1-w \rangle   }\, \frac1{z_3} \ ,
\end{align}
    where $\langle 1-w \rangle $  is the ``average''  value of $1-w$.
    Since $B(w)$  is strongly enhanced for $w=1$, we should expect
    that $\langle 1-w \rangle $ is  numerically small,
    leading to a $\mu \sim k/z_3$ rescaling with a  rather large coefficient $k$.
 As we will see, $k \sim 4$ in this case.


  Again, one may  ask  if
the perturbative formula (\ref{MNL})  involving the
$\ln  z_3^2 $ logarithm may be applied to  actual   lattice data.
 In particular, our exercise with the mass-term IR regularization
 and the  resulting  Bessel function shows
 that
 the logarithmic behavior $\ln z_3^2 $
 of the hard term is valid   only for $z_3 $  values well   below
the IR cut-off $R$, which is given by  the hadron size in our case.
Hence, a practical question is whether
  the data  really
show a logarithmic evolution behavior  in some region of  \mbox{small $z_3$. }
